\section{Model}\label{sec:Model}

The model consists of heterogeneous agents featuring durable goods and directed search. Its objective is to generate business cycle dynamics of durable and nondurable expenditures consistent with empirical patterns.

In short, the timeline of the model is illustrated in figure \ref{fig:timeline}. The agent starts each period with beginning of period durable stock ($n_t$) and market resources ($m_t$). The agent solves the adjuster and keeper problem separately and chooses optimal levels of non-durable consumption $c_t$, durable stock $d_t$ and assets $a_t$. To combine the discrete choices, the agent chooses the optimal search intensity which dictates the success probability of finding a new durable stock. At the end of the period, the agent holds a durable stock $d_t$ and liquid assets $a_t$. Between period $t$ and $t+1$ income shocks materialize, interest rate is received and maintenance costs need to be paid. 

\begin{figure}[h!]
	\centering
	\includegraphics[width=\textwidth]{\FigDir/timeline.pdf}
	\caption{Timeline of the structural model}
	\label{fig:timeline}
\end{figure}

Each household is characterized by two state variables: beginning-of-period market resources $m_t$ and the stock of durables $n_t$. Following \cite{Carroll1992-bh}, all variables are normalized by permanent income and written in lowercase letters. Preferences exhibit constant relative risk aversion and combine nondurable consumption $c_t$ and the service flow of durables $d_t$ using a Cobb-Douglas aggregator\footnote{Using this form is standard (eg see \cite{Berger2015-kp, McKay2021-ib}) which is in line with micro-evidence \cite{Berger2018-yb}. An alternative approach is to use a CES aggregator as in \cite{Mankiw1983-za}}:
\begin{align}
	u(c_t, d_t) &= \frac{\big(c_t^\alpha d_t^{1-\alpha}\big)^{1-\rho}}{1-\rho}, \qquad \alpha \in (0,1), \ \rho > 0.
\end{align}

At the start of the period, the household observes $(m_t, n_t)$ and decides whether to keep the durable stock unchanged or attempt an adjustment. If no adjustment occurs, the household sets $d_t = n_t$ and allocates $m_t$ between nondurable consumption and savings. If the household adjusts, it sells the entire beginning-of-period stock at no price risk but incurs a proportional adjustment cost $\tau \in (0,1)$. The resulting liquidity is
\begin{align}
	\ell_t = m_t + (1-\tau)n_t,
\end{align}
which can then be allocated across nondurable consumption, a new durable stock $d_t$, and savings $a_t$.

Adjustment is not frictionless. To reflect realistic frictions in the durable goods market, the model introduces a directed search mechanism. Households that wish to adjust must exert search effort $\varepsilon_t \in [0,1]$, which determines the probability of finding a suitable durable. The search effort incurs a utility cost scaled by permanent income:
\begin{align}
	\sigma(\varepsilon_t) = S\big[\varepsilon_t + (1-\varepsilon_t)\log(1-\varepsilon_t)\big],
\end{align}
where $S$ is a smoothing parameter that governs the degree of friction in the market. The expected time to find a match is $1/\varepsilon_t$, so a higher $\varepsilon_t$ implies a shorter expected duration between initiating search and obtaining a good match.  

The household's problem before learning whether an adjustment opportunity arrives is:
\begin{align}
	v_t(m_t,n_t) = \max_{\varepsilon_t} \Big[(1-\varepsilon_t) v_t^{\text{keep}}(m_t,n_t) + \varepsilon_t v_t^{\text{adj}}(\ell_t) - \sigma(\varepsilon_t)\Big],
\end{align}
where $v_t^{\text{keep}}$ and $v_t^{\text{adj}}$ are the value functions conditional on keeping or adjusting.

Labor income follows an idiosyncratic process with permanent and transitory components:
\begin{align}
	y_t &=
	\begin{cases}
		(1-T) h w \psi_t \theta_t, & \text{with probability } P(\text{emp}),\\
		y_{\text{unemp}}, & \text{with probability } 1 - P(\text{emp}),
	\end{cases}
\end{align}
where $h$ is hours worked, $w$ is the real wage, and $\psi_t$ and $\theta_t$ are i.i.d.\ log-normally distributed permanent and transitory shocks with $\mathbb{E}[\psi_t] = \mathbb{E}[\theta_t] = 1$. Unemployment occurs with probability $1-P(\text{emp})$, in which case the agent receives unemployment benefits $y_{\text{unemp}}$ expressed as a fraction of permanent income. 

Assets earn an interest rate $r_t$ when positive and $r_t^B$ if negative. The durable stock depreciates at rate $\delta \in (0,1)$ from which a mandatory maintenance cost $\chi$ needs to be paid \cite{Bachmann2013-ka}. Therefore, normalization implies that the next-period states evolve according to:
\begin{align}
	m_{t+1} &= \frac{\mathcal{R} a_t - d_t (\delta \chi)}{\psi_{t+1}} + \theta_{t+1},\\
	n_{t+1} &= \frac{(1-\delta(1 - \chi))d_t}{\psi_{t+1}}.
\end{align}

Where:
	\begin{align}
	\mathcal{R} = 
	\begin{cases}
		1+r_t, & a_t \ge 0,\\
		1+r_t^B, & a_t < 0.
	\end{cases}
\end{align}

Borrowing is limited by a lower bound $\underline{a}$ and a collateral constraint allowing borrowing against the next-period value of durables, up to a loan-to-value ratio $LTV$.

Hence, the Bellman equations conditional on keep and adjust are:
\begin{align}
	v_t^{\text{keep}}(m_t,n_t) &= \max_{c_t} \ u(c_t ,d_t) + \beta \mathbb{E}_t \big[ v_{t+1}(m_{t+1},n_{t+1}) \big],\\
	& \text{s.t. } a_t = m_t - c_t \geq \underline{a} - (1 - \delta(1 - \chi)) (1 - \tau) LTV \cdot d_t,\\
	m_{t+1} &= \frac{\mathcal{R} a_t - d_t (\delta \chi)}{\psi_{t+1}} + \theta_{t+1},\\
	n_{t+1} &= \frac{(1-\delta(1 - \chi))d_t}{\psi_{t+1}}.\\[1em]
	v_t^{\text{adj}}(\ell_t) &= \max_{c_t,d_t} \ u(c_t,d_t) + \beta \mathbb{E}_t \big[ v_{t+1}(m_{t+1},n_{t+1}) \big],\\
	& \text{s.t. } a_t = \ell_t - c_t - d_t \geq \underline{a} - (1 - \delta(1 - \chi)) (1 - \tau) LTV \cdot d_t,\\
	m_{t+1} &= \frac{\mathcal{R} a_t - d_t (\delta \chi)}{\psi_{t+1}} + \theta_{t+1},\\
	n_{t+1} &= \frac{(1-\delta(1 - \chi))d_t}{\psi_{t+1}}.
\end{align}

With the following first order conditions with respect to non-durable consumption and durable stock:

\begin{align}
	\frac{\partial v_t^{adj}}{\partial c_t} &= u_c(c_t, d_t) \geq \beta R \mathbb{E}_t [\frac{v_{t+1}^m}{\psi_{t+1}}], \label{eq:FOC_c} \\
	\frac{\partial v_t^{adj}}{\partial d_t} &= u_d(c_t, d_t) \geq \beta (R + \delta \chi) \mathbb{E}_t[\frac{v_{t+1}^m}{\psi_{t+1}}] -\beta (1 - \delta(1 - \chi)) \mathbb{E}_t[\frac{v_{t+1}^n}{\psi_{t+1}}], \label{eq:FOC_d}
\end{align}

Equations~\ref{eq:FOC_c} and~\ref{eq:FOC_d} are stated as KKT conditions. Equality holds when the collateral constraint does not bind, and the left-hand side strictly exceeds the right-hand side when it does. Equation~\ref{eq:FOC_c} equates the marginal utility of nondurable consumption with the discounted shadow value of savings. Equation~\ref{eq:FOC_d} shows that the marginal utility of durables equals the net user cost, where the first term on the right is the discounted financing cost combining the opportunity cost $\mathcal{R}$ and the maintenance payment $\delta\chi$, and the second term is the discounted benefit of carrying the remaining durable stock $(1-\delta(1-\chi))d_t$ into the next period. 

I describe the solution method in appendix \ref{sec:appdx_model}. In short, I solve the keeper problem using EGM with an upper envelope algorithm and the adjuster using root-finding and FUES to select the optimal points. Smoothing of the policy functions stems from directed search. Since around the indifference points between two options, agents will choose a search intensity between zero and one, the combined policy function is the linear combination of both options. Furthermore, I calculate end of period variables using continuous shock points instead of approximating the income distribution by a few points. 

I solve the model in infinite horizon and calculate Jacobians using the Fake-News algorithm from \cite{Auclert2021-gk} with respect to interest rates, labor income, and the unemployment probability.